\scanpage{48}
termes de (6), les produits fibrés
multiples de $X_{\Delta_1}$ sur $X_{\Delta_0}$, pour $r\in\mathbb N$
\[
(10)\quad\left\{\begin{aligned}
X_{\Delta_r}&=\{(x_0,\ldots,x_r)\in X_{\Delta_0}^{r+1}\mid x_0\leq\cdots x_r\}\\
&\simeq\left\{(x,i_0,\ldots,i_r)\in X\times I^{r+1}\ \middle|\ \begin{matrix}x\in X_{i_0}\\i_0\leq\cdots\leq i_r\end{matrix}\right\}\\
&\simeq\coprod_{i_*\in I(\Delta_r)}\cdot X_{i_*}\\
&\simeq\coprod_{i_0\in I}X_{i_0}\times I_i(\Delta_{r-1})
\end{aligned}\right.
\]
où
\[
(11)\quad\left\{\begin{aligned}
I\unsure{(\Delta_r)|\{\Delta_r\}}&=\operatorname{Hom}_{\text{ens. ord}}(\Delta_r,I)
\quad\begin{matrix}\text{ens des}\\\text{\uline{drapeaux} de}\\\text{longueur }r\text{ de }I\end{matrix}\\
X_{i_*}&=X_{i_0,i_1,\ldots,i_r}\overset{\mathrm{def}}{=}X_{i_0}\quad\text{si }i_*\in I(\Delta_r).
\end{aligned}\right.
\]
Donc on peut considérer, soit
\[
X_{\Delta_r}\subset X_{\Delta_0}^{r+1}
\]
soit
\[
X_{\Delta_r}\subset X\times I^{r+1}.
\]
L'homomorphisme ``de composition'' (9)
permet alors de façon bien connue de
considérer les $X_{\Delta_r}$ comme étant les
composantes d'un espace $\frac12$ simplicial
\[
(12)\qquad X_{\Delta_*}=(X_{\Delta_r})_{r\geq0}
\]
\unsure{où} l'application $\frac12$ simpliciale
\[
(13)\qquad\alpha^*:X_{\Delta_r}\longrightarrow X_{\Delta_{r'}}
\]
associé à une application croissante
\mspage{47}
